\section*{Chapter \ref{ch:g8:binary-and-hexadecimal} --- Binary and Hexadecimal}

\begin{solution}{exo:g8:binary-and-hexadecimal:1}
\code{bin(10)} is \code{"0b1010"} ($8 + 2$), and \code{int("1101", 2)} is
$8 + 4 + 1 = 13$.
\end{solution}

\begin{solution}{exo:g8:binary-and-hexadecimal:2}
\code{hex(255)} is \code{"0xff"}: $255 = 15 \times 16 + 15$. \code{int("1F", 16)}
is $1 \times 16 + 15 = 31$.
\end{solution}

\begin{solution}{exo:g8:binary-and-hexadecimal:3}
The divisions $13 = 2 \times 6 + 1$, $6 = 2 \times 3 + 0$, $3 = 2 \times 1 + 1$
and $1 = 2 \times 0 + 1$, each on its line, then \code{binary: 1101}.
\end{solution}

\begin{solution}{exo:g8:binary-and-hexadecimal:4}
\omcode{code/grade-8/01-binary-and-hexadecimal/exercises/to_binary_fn.py}{1}{12}{Decimal
to \omterm{def:g5:bits-and-bytes:binary}{binary}, as a \omterm{def:g7:functions:function}{function}.}\label{lst:g8:binary-and-hexadecimal:sol-to-binary}

For 0 the \omterm{def:g4:algorithms-repeat-and-decide:loop}{loop} never runs, so the \omterm{def:g7:functions:function}{function} must return \code{"0"} itself; an
empty \omterm{def:g7:strings:string}{string} would not be a number. It prints \code{1101 0 11111111}, and a
test of the chapter compares it with \code{bin} for 0 to 599.
\end{solution}

\begin{solution}{exo:g8:binary-and-hexadecimal:5}
It prints \code{011}: each remainder is added at the end, so the first one,
the \omterm{def:g5:bits-and-bytes:bit}{bit} on the right, comes out first. Adding each \omterm{def:g5:bits-and-bytes:bit}{bit} in front,
\texttt{bits = str(n \% 2) + bits}, gives \code{110}.
\end{solution}

\begin{solution}{exo:g8:binary-and-hexadecimal:6}
2 is not a \omterm{def:g5:bits-and-bytes:binary}{binary} digit, so \code{"102"} is not a \omterm{def:g5:bits-and-bytes:binary}{binary number}, and Python
reports a \code{ValueError}.

\omcode{code/grade-8/01-binary-and-hexadecimal/exercises/is_binary.py}{1}{9}{A
check before converting.}\label{lst:g8:binary-and-hexadecimal:sol-is-binary}

The empty text is refused too: it has no digit at all.
\end{solution}

\begin{solution}{exo:g8:binary-and-hexadecimal:7}
Put zeros in front of the shorter \omterm{def:g7:strings:string}{string} until both have the same length,
then add as before:

\omcode{code/grade-8/01-binary-and-hexadecimal/exercises/add_any.py}{1}{7}{Making
the lengths equal.}\label{lst:g8:binary-and-hexadecimal:sol-add-any}

The rest is \cref{lst:g8:binary-and-hexadecimal:add}; it prints \code{10001}
($11 + 6 = 17$) and \code{10000} ($1 + 15 = 16$).
\end{solution}

\begin{solution}{exo:g8:binary-and-hexadecimal:8}
The first digit is the quotient by 16, the second the remainder:

\omcode{code/grade-8/01-binary-and-hexadecimal/exercises/to_code.py}{1}{13}{Three
numbers to a colour code.}\label{lst:g8:binary-and-hexadecimal:sol-code}

It prints \texttt{\#FF8800 \#2E8B57 \#000000}, the codes of
\cref{out:g8:binary-and-hexadecimal:colours} and black.
\end{solution}

\begin{solution}{exo:g8:binary-and-hexadecimal:9}
$2^{12} = 4\,096$ patterns; the largest number is $4\,095$.
\end{solution}

\begin{solution}{exo:g8:binary-and-hexadecimal:10}
$2^{16} - 1 = 65\,535$, written \code{FFFF}: sixteen ones make four groups of
\code{1111}, four \code{F}s.
\end{solution}

\begin{solution}{exo:g8:binary-and-hexadecimal:11}
Every number from 0 to $2^n - 1$ has a writing with $n$ \omterm{def:g5:bits-and-bytes:bit}{bits}: repeated
division gives its \omterm{def:g5:bits-and-bytes:bit}{bits}, at most $n$ of them since the number is below
$2^n$, and zeros in front fill the $n$ places. So each of the $2^n$ numbers
gets at least one of the patterns. There are exactly $2^n$ patterns, and each
pattern is the writing of only one number. If one number had two patterns,
the $2^n$ numbers would use more than $2^n$ patterns, which do not exist.
So each number has exactly one.
\end{solution}

\begin{solution}{exo:g8:binary-and-hexadecimal:12}
Read the text into a whole number (the doubling walk of
\cref{lst:g8:binary-and-hexadecimal:to-decimal}, with the base instead of 2),
then write that number in the other base (repeated division):

\omcode{code/grade-8/01-binary-and-hexadecimal/exercises/convert.py}{1}{26}{Any
base to any base.}\label{lst:g8:binary-and-hexadecimal:sol-convert}

It prints \code{2C 11111111} and \code{377 200221}. A test converts many
numbers to every base from 2 to 16 and back.
\end{solution}

\begin{solution}{pb:g8:binary-and-hexadecimal:1}
\textbf{1.} 1, 2, 4, 8, 16, 32, 64, 128: square $k$ holds $2^{k-1}$ grains.

\textbf{2.} $1 + 2 + \dots + 128 = 255$, one less than the 256 grains of
square 9.

\textbf{3.} Square $k$ holds a 1 followed by $k - 1$ zeros; the total on 8
squares is \code{11111111}: each square adds a different power of 2, so each
place has exactly one 1.

\textbf{4.} The first square must count 1 grain; doubling first would start
the board at 2.

\textbf{5.} The \omterm{def:g4:programs-and-operating-system:program}{program} prints:

\omoutput{code/grade-8/01-binary-and-hexadecimal/outputs/grains.txt}{The rice on
the board.}\label{out:g8:binary-and-hexadecimal:grains}

The last square holds $2^{63}$ grains, more than nine billion billion, and
the total is almost twice as much.

\textbf{6.} The total is $1 + 2 + \dots + 2^{63}$: in \omterm{def:g5:bits-and-bytes:binary}{binary}, 64 ones. By the
proposition, the largest number on 64 \omterm{def:g5:bits-and-bytes:bit}{bits} is $2^{64} - 1$; the \omterm{def:g4:programs-and-operating-system:program}{program}'s two
lines agree.

\textbf{7.} 64 ones, so 16 groups of \code{1111}: sixteen \code{F}s.

\textbf{8.} About $1.8 \times 10^{19}$ grains of 25 milligrams make about
$4.6 \times 10^{20}$ milligrams, that is about $4.6 \times 10^{11}$ tonnes:
some 460 billion tonnes of rice.

\textbf{9.} 64 \omterm{def:g5:bits-and-bytes:bit}{bits}, which is 8 \omterm{def:g4:memory-and-storage:byte}{bytes}.

\textbf{10.} 32 \omterm{def:g5:bits-and-bytes:bit}{bits} count up to $2^{32} - 1 = 4\,294\,967\,295$, which is
exactly the total on the first 32 squares; the 33rd square goes past it.

\textbf{11.} 64 \omterm{def:g5:bits-and-bytes:bit}{bits}: the total is the largest number they can write.

\textbf{12.} \textbf{Result.} The board holds 18\,446\,744\,073\,709\,551\,615
grains, $2^{64} - 1$, written as sixteen \code{F}s in \omterm{def:g8:binary-and-hexadecimal:hexadecimal}{hexadecimal}: exactly the
largest number that 8 \omterm{def:g4:memory-and-storage:byte}{bytes} can hold.
\end{solution}
